% test-002.tex -- comandos \spunit, \spnewunit y \spunitalias
%
% Compilar:  lualatex-dev test-002.tex
% Validar:   show-pdf-tags --xml test-002.pdf | rnv-wrapp
%            verapdf --flavour ua2 --format text test-002.pdf
%
% Cada caso es un parrafo numerado, para poder referirse a el al
% escuchar el PDF. Los comentarios "doc:" indican la lectura que
% documenta el manual de usuario; en el resto, la lectura se revisa
% escuchando. Las entradas invalidas (que detienen la compilacion)
% no van aqui. La tabla completa de unidades y alias va en test-003.
\DocumentMetadata
  {
    lang=es-CL,
    pdfversion=2.0,
    pdfstandard=ua-2,
    tagging=on,
    tagging-setup={math/setup={mathml-SE}},
  }
\documentclass{article}
\usepackage[spanish]{babel}
\usepackage{unicode-math}
\usepackage{spintent}
\usepackage{hyperref}
\hypersetup
  {
    pdfdisplaydoctitle = true,
    pdftitle           = {test-002: spunit},
  }
\begin{document}

\section*{A. Formas del argumento}

% doc: «metro»
1. Solo la unidad: $\spunit{m}$.

% doc: «23 metros»
2. Número y unidad: $\spunit{23 m}$.

% doc: «23 metros segundos cuadrado»
3. Producto de unidades con potencia: $\spunit{23 m*s^{2}}$.

% doc: «23 metros por segundos cuadrado»
4. Fracción lineal: $\spunit{23 m/s^{2}}$.

% doc: «2 newtons metros»
5. Producto de dos unidades: $\spunit{2 N*m}$.

% doc: «3 kilogramos metros segundos»
6. Producto de tres unidades: $\spunit{3 kg*m*s}$.

% doc: «3 kilogramos metros cuadrado por segundos cuadrado»
7. Producto y fracción: $\spunit{3 kg*m^{2}/s^{2}}$.

% doc: «9 coma 8 metros por segundos cuadrado»
8. Número decimal: $\spunit{9,8 m/s^{2}}$.

% doc: «menos 3 coma 5 metros»
9. Número negativo: $\spunit{-3,5 m}$.

% doc: «1 coma 2 con 3 periódico metros»
10. Número periódico: $\spunit{1,2;3 m}$.

% doc: «3 por 10 a la octavo potencia metros por segundos»
11. Notación científica: $\spunit{3e8 m/s}$.

% doc: «metro cuadrado»
12. Sin número, con potencia: $\spunit{m^{2}}$.

% doc: «metros por segundos»
13. Sin número, con fracción: $\spunit{m/s}$.

% doc: «23 metros»
14. Con alias de la tabla: $\spunit{23 metros}$.

% doc: «20 grados Celsius»
15. Con alias de temperatura: $\spunit{20 celsius}$.

\section*{B. Llave cifras-sep}

% doc: «100000 metros»
16. Por defecto: $\spunit{100000 m}$.

% doc: «100000 metros»
17. Con false y espacios del usuario: $\spunit[cifras-sep=false]{10\,00\,00 m}$.

% doc: «100000 metros»
18. Con false y sin espacios: $\spunit[cifras-sep=false]{100000 m}$.

\section*{C. Llave read-period-sep}

% doc: «3 coma 4 periódico metros»
19. Con coma: $\spunit[read-period-sep=coma]{3;4 m}$.

% doc: «3 punto 4 periódico metros»
20. Con punto: $\spunit[read-period-sep=punto]{3;4 m}$.

\section*{D. Llave notation-sym}

% doc: «3 por 10 a la octavo potencia metros por segundos»
21. Con cdot: $\spunit[notation-sym=cdot]{3e8 m/s}$.

% doc: «3 por 10 a la octavo potencia metros por segundos»
22. Con times: $\spunit[notation-sym=times]{3e8 m/s}$.

\section*{E. Llave print-cdot}

% doc: «23 metros segundos cuadrado»
23. Por defecto: $\spunit{23 m*s^{2}}$.

% doc: «23 metros segundos cuadrado»
24. Con false: $\spunit[print-cdot=false]{23 m*s^{2}}$.

% doc: «23 metros segundos cuadrado»
25. Con true: $\spunit[print-cdot=true]{23 m*s^{2}}$.

% doc: «2 newtons metros»
26. Con true y dos unidades: $\spunit[print-cdot=true]{2 N*m}$.

\section*{F. Llave read-cdot}

% doc: «23 metros segundos cuadrado»
27. Por defecto: $\spunit{23 m*s^{2}}$.

% doc: «23 metros segundos cuadrado»
28. Con false: $\spunit[read-cdot=false]{23 m*s^{2}}$.

% doc: «23 metros segundos cuadrado»
29. Con true: $\spunit[read-cdot=true]{23 m*s^{2}}$.

% doc: «2 newtons por metros»
30. Con true y print-cdot true: $\spunit[print-cdot=true,read-cdot=true]{2 N*m}$.

\section*{G. Llave read-slash}

% doc: «23 metros por segundos cuadrado»
31. Por defecto: $\spunit{23 m/s^{2}}$.

% doc: «23 metros por segundos cuadrado»
32. Con por: $\spunit[read-slash=por]{23 m/s^{2}}$.

% doc: «23 metros partido por segundos cuadrado»
33. Con partido-por: $\spunit[read-slash=partido-por]{23 m/s^{2}}$.

% doc: «metros partido por segundos»
34. Con partido-por y sin número: $\spunit[read-slash=partido-por]{m/s}$.

\section*{H. Combinaciones de llaves}

% doc: «9 coma 8 kilogramos por metros partido por segundos cuadrado»
35. Todas las de unidades: $\spunit[print-cdot=true,read-cdot=true,read-slash=partido-por]{9,8 kg*m/s^{2}}$.

% doc: «3 por 10 a la octavo potencia metros partido por segundos»
36. Llaves de número y de unidades: $\spunit[cifras-sep=false,notation-sym=times,read-slash=partido-por]{3e8 m/s}$.

\section*{I. Unidades y alias nuevos}

\spnewunit{px}{píxeles}%
\spnewunit{ppm}{partes por millón}%
\spunitalias{caudal}{m^{3}/s}%
\spunitalias{velocidad}{m/s}%

% doc: «100 píxeles»
37. Unidad nueva: $\spunit{100 px}$.

% doc: «5 partes por millón»
38. Otra unidad nueva: $\spunit{5 ppm}$.

% doc: «3 píxeles segundos»
39. Unidad nueva en un producto: $\spunit{3 px*s}$.

% doc: «50 metros al cubo por segundos»
40. Alias nuevo: $\spunit{50 caudal}$.

% doc: «30 metros por segundos»
41. Otro alias nuevo: $\spunit{30 velocidad}$.

% doc: «50 metros al cubo partido por segundos»
42. Alias nuevo con read-slash: $\spunit[read-slash=partido-por]{50 caudal}$.

\section*{J. Dentro de fórmulas}

% doc: «3 metros más 2 metros es igual a 5 metros»
43. Suma e igualdad: $\spunit{3 m} + \spunit{2 m} = \spunit{5 m}$.

% doc: «g es igual a 9 coma 8 metros por segundos cuadrado»
44. Con una variable: $g = \spunit{9,8 m/s^{2}}$.

% doc: «3 coma 5 punto 2 metros»
45. Con un número aparte: $\spnum{3,5} \cdot \spunit{2 m}$.

% doc: «2 metros por segundos es menor que 3 metros partido por segundos»
46. Llaves distintas en una misma fórmula: $\spunit[read-slash=por]{2 m/s} < \spunit[read-slash=partido-por]{3 m/s}$.

% doc: «3 kilogramos metros cuadrado por segundos cuadrado»
47. En fórmula destacada:
\[ \spunit{3 kg*m^{2}/s^{2}} \]

% doc: «la fracción con numerador 10 metros y denominador 2 segundos»
48. En una fracción: $\frac{\spunit{10 m}}{\spunit{2 s}}$.

\end{document}
